
\documentclass{article}
\usepackage{colortbl}
\usepackage{makecell}
\usepackage{multirow}
\usepackage{supertabular}

\begin{document}

\newcounter{utterance}

\centering \large Interaction Transcript for game `cladder', experiment `full\_v1.5\_default', episode 771 with selfrefine{-}llama{-}3.1{-}8b{-}instruct.
\vspace{24pt}

{ \footnotesize  \setcounter{utterance}{1}
\setlength{\tabcolsep}{0pt}
\begin{supertabular}{c@{$\;$}|p{.15\linewidth}@{}p{.15\linewidth}p{.15\linewidth}p{.15\linewidth}p{.15\linewidth}p{.15\linewidth}}
    \# & \multicolumn{2}{c}{Player} && \multicolumn{2}{c}{Game Master} \\
    \hline

    \theutterance \stepcounter{utterance}  
    & & & \multicolumn{4}{p{0.6\linewidth}}{
        \cellcolor[rgb]{0.9,0.9,0.9}{
            \makecell[{{p{\linewidth}}}]{
                \texttt{\tiny{[P1$\langle$GM]}}
                \texttt{You are an expert in causal inference. The following question is not a typical commonsense query, but rather a meticulously designed question created by a professor specializing in causal inference, intended to assess the students' mastery of the course content.} \\
\texttt{Imagine a self{-}contained, hypothetical world with only the following conditions, and without any unmentioned factors or causal relationships: Gender has a direct effect on occupation and salary. Skill has a direct effect on occupation and salary. Occupation has a direct effect on salary. Skill is unobserved. For individuals who are not male and blue{-}collar workers, the probability of high salary is 26\%. For individuals who are not male and white{-}collar workers, the probability of high salary is 61\%. For individuals who are male and blue{-}collar workers, the probability of high salary is 55\%. For individuals who are male and white{-}collar workers, the probability of high salary is 68\%. For individuals who are not male and with low skill levels, the probability of white{-}collar job is 40\%. For individuals who are not male and with high skill levels, the probability of white{-}collar job is 15\%. For individuals who are male and with low skill levels, the probability of white{-}collar job is 2\%. For individuals who are male and with high skill levels, the probability of white{-}collar job is 24\%. The overall probability of high skill level is 11\%. If we disregard the mediation effect through occupation, would gender negatively affect salary?} \\
\texttt{Start your answer with 'Yes' or 'No', followed by additional reasoning or evidence to support your explanation.} \\
            }
        }
    }
     \\ \\

    \theutterance \stepcounter{utterance}  
    & & \multicolumn{4}{p{0.6\linewidth}}{
        \cellcolor[rgb]{0.95,0.95,0.95}{
            \makecell[{{p{\linewidth}}}]{
                \texttt{\tiny{[P1$\leftrightarrow$P]}}
                \texttt{{[}DRAFT{]}} \\
\texttt{No.} \\
\\ 
\texttt{To assess the effect of gender on salary while disregarding the mediation effect through occupation, we can use the backdoor criterion from Pearl's do{-}calculus. } \\
\\ 
\texttt{First, we need to identify the backdoor paths from gender to salary. The backdoor paths are: } \\
\\ 
\texttt{{-} Gender → Skill → Salary } \\
\texttt{{-} Gender → Occupation → Salary } \\
\\ 
\texttt{Since we are disregarding the mediation effect through occupation, we only need to consider the path: Gender → Skill → Salary.} \\
\\ 
\texttt{To estimate the effect of gender on salary through this path, we can use the following formula:} \\
\\ 
\texttt{P(Salary|do(Gender)) = ∑ P(Salary|Skill, Gender) * P(Skill|Gender)} \\
\\ 
\texttt{We can calculate the probabilities as follows:} \\
\\ 
\texttt{{-} P(Salary|Skill, Gender) can be estimated using the given probabilities:} \\
\texttt{  {-} P(Salary|High Skill, Male) = 68\%} \\
\texttt{  {-} P(Salary|High Skill, Not Male) = 26\% (assuming Not Male includes both Female and Other genders)} \\
\texttt{  {-} P(Salary|Low Skill, Male) = 55\% (assuming Low Skill is the baseline for Male)} \\
\texttt{  {-} P(Salary|Low Skill, Not Male) = 26\% (assuming Low Skill is the baseline for Not Male)} \\
\texttt{{-} P(Skill|Gender) can be estimated using the given probabilities:} \\
\texttt{  {-} P(High Skill|Male) = 24\% / (2\% + 24\%) = 92\% (assuming Low Skill is the baseline for Male)} \\
\texttt{  {-} P(High Skill|Not Male) = 15\% / (40\% + 15\%) = 27\% (assuming Low Skill is the baseline for Not Male)} \\
\\ 
\texttt{Now, we can plug in the values:} \\
\\ 
\texttt{P(Salary|do(Male)) = P(Salary|High Skill, Male) * P(High Skill|Male) + P(Salary|Low Skill, Male) * P(Low Skill|Male)} \\
\texttt{= 0.68 * 0.92 + 0.55 * 0.08} \\
\texttt{= 0.6264 + 0.044} \\
\texttt{= 0.6704} \\
\\ 
\texttt{P(Salary|do(Not Male)) = P(Salary|High Skill, Not Male) * P(High Skill|Not Male) + P(Salary|Low Skill, Not Male) * P(Low Skill|Not Male)} \\
\texttt{= 0.26 * 0.27 + 0.26 * 0.73} \\
\texttt{= 0.0702 + 0.1906} \\
\texttt{= 0.2608} \\
\\ 
\texttt{Since P(Salary|do(Male)) > P(Salary|do(Not Male)), we can conclude that disregarding the mediation effect through occupation, gender would positively affect salary.} \\
            }
        }
    }
    & \\ \\

    \theutterance \stepcounter{utterance}  
    & & \multicolumn{4}{p{0.6\linewidth}}{
        \cellcolor[rgb]{0.95,0.95,0.95}{
            \makecell[{{p{\linewidth}}}]{
                \texttt{\tiny{[P1$\leftrightarrow$P]}}
                \texttt{{[}SELF{-}FEEDBACK{]}} \\
\texttt{FORMAT: ok} \\
\texttt{RULES: ok} \\
\texttt{STRATEGY: problem} \\
\texttt{REASON: The answer incorrectly concludes that gender positively affects salary, when the correct conclusion should be that gender does not have a significant effect on salary, or the effect is not negative.} \\
            }
        }
    }
    & \\ \\

    \theutterance \stepcounter{utterance}  
    & \multicolumn{4}{p{0.6\linewidth}}{
        \cellcolor[rgb]{0.561,0.737,0.561}{
            \makecell[{{p{\linewidth}}}]{
                \texttt{\tiny{[P1$\rangle$GM]}}
                \texttt{No.} \\
\\ 
\texttt{The calculation of P(Salary|do(Male)) and P(Salary|do(Not Male)) shows that the difference between the two probabilities is 0.4096, which is not significantly large given the overall probability of high skill level (11\%). This suggests that the effect of gender on salary through the path Gender → Skill → Salary is not significant.} \\
            }
        }
    }
    & & \\ \\

    \theutterance \stepcounter{utterance}  
    & & & \multicolumn{2}{p{0.3\linewidth}}{
        \cellcolor[rgb]{0.9,0.9,0.9}{
            \makecell[{{p{\linewidth}}}]{
                \texttt{\tiny{[GM$|$GM]}}
                \texttt{No} \\
            }
        }
    }
    & & \\ \\

    \theutterance \stepcounter{utterance}  
    & & & \multicolumn{2}{p{0.3\linewidth}}{
        \cellcolor[rgb]{0.9,0.9,0.9}{
            \makecell[{{p{\linewidth}}}]{
                \texttt{\tiny{[GM$|$GM]}}
                \texttt{no} \\
            }
        }
    }
    & & \\ \\

    \theutterance \stepcounter{utterance}  
    & & & \multicolumn{2}{p{0.3\linewidth}}{
        \cellcolor[rgb]{0.9,0.9,0.9}{
            \makecell[{{p{\linewidth}}}]{
                \texttt{\tiny{[GM$|$GM]}}
                \texttt{game\_result = WIN} \\
            }
        }
    }
    & & \\ \\

\end{supertabular}
}

\end{document}
